% ============================================================
%  GUÍA 1: OPERACIONES EN Q (FRACCIONES, POTENCIACIÓN Y ECUACIONES)
%  Curso de Introducción 2do Año
%  Autor: Ing. Gabriel Astudillo
%  Clases cubiertas: 1 a 5
% ============================================================

\documentclass[12pt, a4paper]{article}

% ── Paquetes ─────────────────────────────────────────────────

\usepackage{mathwizards-guia}
\mwguia{Guía 1 — Operaciones en $\mathbb{Q}$}{Ing. Gabriel Astudillo $\cdot$ 2do Año}

% ── Comandos propios de esta guía ───────────────────
\newcommand{\Q}{\mathbb{Q}}
\newcommand{\Z}{\mathbb{Z}}

\begin{document}
% ============================================================

% ── Portada / Título ─────────────────────────────────────────
\mwportada{Guía de Estudio 1}{Operaciones en $\mathbb{Q}$}{Fracciones, Potenciación y Ecuaciones con Fracciones}

\vspace{4pt}
\begin{cuadroinfo}
  \small
  \textbf{Presentación:} Los números racionales ($\Q$) están en todas partes: en la cocina, en la construcción, en las mediciones y en la vida diaria. En esta guía aprenderás a operar con fracciones con seguridad: sumar, restar, multiplicar y dividir, resolver operaciones combinadas, aplicar las leyes de los exponentes a expresiones fraccionarias y plantear y resolver ecuaciones lineales con fracciones. Cada sección incluye \textbf{ejercicios resueltos} y \textbf{ejercicios propuestos} para que practiques.
  \\[4pt]
  \textbf{Nivel de Dificultad:}\quad
  \facil{} Básico / Directo \quad
  \medio{} Intermedio \quad
  \dificil{} Desafiante
\end{cuadroinfo}

\vspace{8pt}

% ============================================================
\section{Fracciones: Definición y Tipos}
% ============================================================

\subsection{¿Qué es una fracción?}

\begin{cuadroteoria}
  \textbf{Fracción:} es la representación de una parte de un todo o de una división entre dos números. Se escribe
  $$\frac{a}{b} \qquad \text{con } b \neq 0,$$
  donde:
  \begin{itemize}[leftmargin=*]
    \item $a$ es el \textbf{numerador}: indica cuántas partes se toman.
    \item $b$ es el \textbf{denominador}: indica en cuántas partes iguales se divide el todo.
  \end{itemize}

  \textbf{Tipos de fracciones:}
  \begin{itemize}[leftmargin=*]
    \item \textbf{Propia:} el numerador es menor que el denominador ($\frac{3}{5}$). Representa menos de un entero.
    \item \textbf{Impropia:} el numerador es mayor o igual que el denominador ($\frac{7}{4}$). Representa más de un entero.
    \item \textbf{Mixta:} combina un entero y una fracción propia ($2\frac{1}{3}$). Es equivalente a una fracción impropia.
    \item \textbf{Equivalentes:} representan la misma cantidad pero con diferentes números ($\frac{1}{2} = \frac{2}{4} = \frac{3}{6}$).
  \end{itemize}
\end{cuadroteoria}

\begin{cuadropasos}
  \textbf{Convertir mixta a impropia:} $a\frac{b}{c} = \dfrac{a \cdot c + b}{c}$ \quad (se multiplica el entero por el denominador y se suma el numerador).

  \textbf{Convertir impropia a mixta:} se divide el numerador entre el denominador; el cociente es el entero, el resto es el nuevo numerador y el divisor se conserva.
\end{cuadropasos}

\begin{cuadroteoria}
  \textbf{Simplificación de fracciones:} una fracción se simplifica dividiendo numerador y denominador por un divisor común hasta llegar a la \textbf{fracción irreductible} (aquella cuyo numerador y denominador no tienen divisores comunes distintos de 1).
\end{cuadroteoria}

\subsection{Ejercicios resueltos}

\textbf{Ejemplo R1.} Clasifica las siguientes fracciones en propias o impropias y escribe las impropias como número mixto.

$$\frac{2}{7}, \qquad \frac{9}{4}, \qquad \frac{5}{8}, \qquad \frac{11}{3}$$

\begin{cuadrosolucion}
  \textbf{Solución:}
  \begin{itemize}
    \item $\dfrac{2}{7}$: propia (el numerador $2 < 7$).
    \item $\dfrac{9}{4}$: impropia. Como $9 \div 4$ da cociente $2$ y resto $1$: $\dfrac{9}{4} = 2\frac{1}{4}$.
    \item $\dfrac{5}{8}$: propia ($5 < 8$).
    \item $\dfrac{11}{3}$: impropia. $11 \div 3$ da cociente $3$ y resto $2$: $\dfrac{11}{3} = 3\frac{2}{3}$.
  \end{itemize}
\end{cuadrosolucion}

\textbf{Ejemplo R2.} Escribe como fracción impropia: $a) \; 3\frac{2}{5} \qquad b) \; 7\frac{1}{2}$

\begin{cuadrosolucion}
  \textbf{Solución:}
  \begin{itemize}
    \item $a) \; 3\frac{2}{5} = \dfrac{3 \cdot 5 + 2}{5} = \dfrac{15 + 2}{5} = \dfrac{17}{5}$.
    \item $b) \; 7\frac{1}{2} = \dfrac{7 \cdot 2 + 1}{2} = \dfrac{14 + 1}{2} = \dfrac{15}{2}$.
  \end{itemize}
\end{cuadrosolucion}

\newpage %ajuste para pág

\textbf{Ejemplo R3.} Simplifica hasta la fracción irreductible: $a) \; \dfrac{12}{18} \qquad b) \; \dfrac{45}{60}$

\begin{cuadrosolucion}
  \textbf{Solución:}
  \begin{itemize}
    \item $a) \; \dfrac{12}{18} = \dfrac{12 \div 6}{18 \div 6} = \dfrac{2}{3}$.
    \item $b) \; \dfrac{45}{60} = \dfrac{45 \div 15}{60 \div 15} = \dfrac{3}{4}$.
  \end{itemize}
\end{cuadrosolucion}

\textbf{Ejemplo R4.} Completa para que las fracciones sean equivalentes: $\dfrac{3}{4} = \dfrac{?}{20}$

\begin{cuadrosolucion}
  \textbf{Solución:} Como $20 \div 4 = 5$, multiplicamos numerador y denominador por $5$:
  $$\frac{3}{4} = \frac{3 \cdot 5}{4 \cdot 5} = \frac{15}{20}.$$
  La fracción buscada es $\dfrac{15}{20}$.
\end{cuadrosolucion}

%\newpage

\subsection{Ejercicios propuestos}

\begin{enumerate}[leftmargin=*]
  \item \facil{} Clasifica en propia o impropia: $a) \; \dfrac{5}{9} \qquad b) \; \dfrac{7}{3} \qquad c) \; \dfrac{11}{11} \qquad d) \; \dfrac{3}{10}$

  \item \facil{} Escribe como fracción impropia: $a) \; 2\frac{3}{4} \qquad b) \; 5\frac{1}{6} \qquad c) \; 8\frac{2}{3}$

  \item \facil{} Escribe como número mixto: $a) \; \dfrac{17}{5} \qquad b) \; \dfrac{22}{7} \qquad c) \; \dfrac{31}{4}$

  \item \facil{} Simplifica hasta la fracción irreductible: $a) \; \dfrac{8}{12} \qquad b) \; \dfrac{24}{36} \qquad c) \; \dfrac{50}{75}$

  \item \facil{} Completa para obtener fracciones equivalentes: $a) \; \dfrac{2}{5} = \dfrac{?}{15} \qquad b) \; \dfrac{7}{8} = \dfrac{?}{32} \qquad c) \; \dfrac{5}{6} = \dfrac{20}{?}$

  \item \medio{} Una pizza se divide en 8 porciones iguales. Comí 3 porciones. ¿Qué fracción de la pizza me queda? ¿Es propia o impropia?

  \item \medio{} Un atleta corre $\dfrac{3}{4}$ de kilómetro cada día. ¿Qué distancia recorre en 2 días? Expresa el resultado como número mixto y como fracción impropia.

  \item \dificil{} Ordena de menor a mayor: $\dfrac{3}{4}, \; \dfrac{5}{6}, \; \dfrac{7}{12}$ (pista: lleva todas a denominador común 12).
\end{enumerate}

\newpage

% ============================================================
\section{Suma y Resta de Fracciones}
% ============================================================

\subsection{Con igual denominador}

\begin{cuadroteoria}
  \textbf{Suma y resta con igual denominador:} se suman (o restan) los numeradores y se conserva el denominador.
  $$\frac{a}{c} + \frac{b}{c} = \frac{a + b}{c}, \qquad \frac{a}{c} - \frac{b}{c} = \frac{a - b}{c}$$
  Siempre que sea posible, se simplifica el resultado.
\end{cuadroteoria}

\subsection{Con diferente denominador}

\begin{cuadroteoria}
  \textbf{Método del producto cruzado:} para sumar o restar $\dfrac{a}{b} \pm \dfrac{c}{d}$:
  $$\frac{a}{b} \pm \frac{c}{d} = \frac{a \cdot d \pm c \cdot b}{b \cdot d}.$$
  Se multiplica \emph{en cruz} y el denominador es el producto de ambos denominadores.
\end{cuadroteoria}

\begin{cuadroteoria}
  \textbf{Método del mínimo común múltiplo (mcm):}
  \begin{enumerate}[leftmargin=*]
    \item Se calcula el mcm de los denominadores (mínimo común múltiplo).
    \item Se divide el mcm entre cada denominador y se multiplica por el numerador correspondiente.
    \item Se suman o restan los nuevos numeradores sobre el mcm.
    \item Se simplifica si es posible.
  \end{enumerate}
\end{cuadroteoria}

\begin{cuadroadvertencia}
  \textbf{Errores clásicos:}
  \begin{itemize}[leftmargin=*]
    \item Sumar numeradores \emph{y} denominadores: $\dfrac{1}{2} + \dfrac{1}{3} \neq \dfrac{2}{5}$. El resultado correcto es $\dfrac{5}{6}$.
    \item Olvidar simplificar la respuesta final.
    \item Confundir la resta: al restar, el signo afecta a \emph{todo} el numerador del sustraendo.
  \end{itemize}
\end{cuadroadvertencia}

\newpage %ajuste para pág

\subsection{Ejercicios resueltos}

\textbf{Ejemplo R1.} Calcula: $a) \; \dfrac{2}{7} + \dfrac{3}{7} \qquad b) \; \dfrac{5}{9} - \dfrac{2}{9}$

\begin{cuadrosolucion}
  \textbf{Solución:}
  \begin{itemize}
    \item $a) \; \dfrac{2}{7} + \dfrac{3}{7} = \dfrac{2 + 3}{7} = \dfrac{5}{7}$.
    \item $b) \; \dfrac{5}{9} - \dfrac{2}{9} = \dfrac{5 - 2}{9} = \dfrac{3}{9} = \dfrac{1}{3}$ (simplificando por 3).
  \end{itemize}
\end{cuadrosolucion}

\textbf{Ejemplo R2.} Calcula por el método del producto cruzado: $\dfrac{2}{3} + \dfrac{1}{4}$

\begin{cuadrosolucion}
  \textbf{Solución:}
  $$\frac{2}{3} + \frac{1}{4} = \frac{2 \cdot 4 + 1 \cdot 3}{3 \cdot 4} = \frac{8 + 3}{12} = \frac{11}{12}.$$
\end{cuadrosolucion}

\textbf{Ejemplo R3.} Calcula por el método del mcm: $\dfrac{5}{6} - \dfrac{3}{4}$

\begin{cuadrosolucion}
  \textbf{Solución:}
  \begin{enumerate}
    \item mcm$(6, 4) = 12$.
    \item $\dfrac{5}{6} = \dfrac{5 \cdot 2}{12} = \dfrac{10}{12}$ \quad y \quad $\dfrac{3}{4} = \dfrac{3 \cdot 3}{12} = \dfrac{9}{12}$.
    \item Restamos: $\dfrac{10}{12} - \dfrac{9}{12} = \dfrac{1}{12}$.
  \end{enumerate}
\end{cuadrosolucion}

\textbf{Ejemplo R4.} Calcula: $2\frac{1}{2} + \dfrac{3}{4}$

\begin{cuadrosolucion}
  \textbf{Solución:} Convertimos el número mixto a fracción impropia:
  $$2\frac{1}{2} = \frac{5}{2}.$$
  Luego: $\dfrac{5}{2} + \dfrac{3}{4}$. Como mcm$(2,4) = 4$:
  $$\frac{5}{2} + \frac{3}{4} = \frac{10}{4} + \frac{3}{4} = \frac{13}{4} = 3\frac{1}{4}.$$
\end{cuadrosolucion}

\newpage

\subsection{Ejercicios propuestos}

\begin{enumerate}[leftmargin=*, resume]
  \item \facil{} Calcula: $a) \; \dfrac{1}{8} + \dfrac{3}{8} \qquad b) \; \dfrac{7}{10} - \dfrac{2}{10} \qquad c) \; \dfrac{4}{9} + \dfrac{5}{9}$

  \item \facil{} Calcula (producto cruzado): $a) \; \dfrac{1}{2} + \dfrac{1}{3} \qquad b) \; \dfrac{2}{5} - \dfrac{1}{4} \qquad c) \; \dfrac{3}{7} + \dfrac{2}{5}$

  \item \facil{} Calcula (mcm): $a) \; \dfrac{1}{6} + \dfrac{3}{4} \qquad b) \; \dfrac{5}{8} - \dfrac{1}{3} \qquad c) \; \dfrac{2}{5} + \dfrac{7}{10}$

  \item \medio{} Calcula: $a) \; 1\frac{1}{3} + \dfrac{5}{6} \qquad b) \; 3\frac{1}{4} - 1\frac{1}{2}$

  \item \medio{} Calcula: $a) \; \dfrac{3}{4} + \dfrac{5}{6} - \dfrac{1}{3} \qquad b) \; \dfrac{7}{12} - \dfrac{1}{4} + \dfrac{2}{3}$

  \item \medio{} Para una receta necesito $\dfrac{2}{3}$ de taza de azúcar y $\dfrac{1}{4}$ de taza de harina. ¿Cuántas tazas de ingredientes uso en total?

  \item \dificil{} Entre tres amigos se reparten una herencia: el primero recibe $\dfrac{1}{3}$, el segundo $\dfrac{1}{4}$ y el tercero el resto. ¿Qué fracción del total recibe el tercero?
\end{enumerate}

%\newpage %ajuste pag

% ============================================================
\section{Multiplicación y División de Fracciones}
% ============================================================

\subsection{Multiplicación}

\begin{cuadroteoria}
  \textbf{Multiplicación de fracciones:} se multiplican los numeradores entre sí y los denominadores entre sí.
  $$\frac{a}{b} \cdot \frac{c}{d} = \frac{a \cdot c}{b \cdot d}$$

  \textbf{Consejo:} si es posible, \emph{simplifica antes de multiplicar} (cancelando factores comunes entre numeradores y denominadores).
\end{cuadroteoria}

\subsection{División}

\begin{cuadroteoria}
  \textbf{División de fracciones:} se multiplica la primera fracción por el \textbf{recíproco} (inverso multiplicativo) de la segunda.
  $$\frac{a}{b} \div \frac{c}{d} = \frac{a}{b} \cdot \frac{d}{c} = \frac{a \cdot d}{b \cdot c}$$

  Alternativamente, se puede aplicar la regla del \emph{producto cruzado}: $\dfrac{a}{b} \div \dfrac{c}{d} = \dfrac{a \cdot d}{b \cdot c}$.
\end{cuadroteoria}

\newpage %ajuste para pág

\begin{cuadroadvertencia}
  \textbf{Cuidado:}
  \begin{itemize}[leftmargin=*]
    \item Para dividir se \emph{invierte} la segunda fracción, no la primera.
    \item Un entero se puede escribir como fracción con denominador 1: $3 = \dfrac{3}{1}$.
    \item $a \div \dfrac{b}{c} = a \cdot \dfrac{c}{b}$.
  \end{itemize}
\end{cuadroadvertencia}

\subsection{Ejercicios resueltos}

\textbf{Ejemplo R1.} Calcula: $a) \; \dfrac{2}{3} \cdot \dfrac{4}{5} \qquad b) \; \dfrac{3}{8} \cdot \dfrac{4}{9}$

\begin{cuadrosolucion}
  \textbf{Solución:}
  \begin{itemize}
    \item $a) \; \dfrac{2}{3} \cdot \dfrac{4}{5} = \dfrac{2 \cdot 4}{3 \cdot 5} = \dfrac{8}{15}$.
    \item $b) \;$ Simplificamos antes de multiplicar: $\dfrac{3}{8} \cdot \dfrac{4}{9}$. El 3 y el 9 se simplifican por 3, y el 8 y el 4 por 4:
          $$\frac{\cancel{3}}{8} \cdot \frac{4}{\cancel{9}} = \frac{1}{\cancel{8}^{2}} \cdot \frac{\cancel{4}^{1}}{3} = \frac{1}{6}.$$
  \end{itemize}
\end{cuadrosolucion}

\textbf{Ejemplo R2.} Calcula: $\dfrac{5}{6} \div \dfrac{2}{3}$

\begin{cuadrosolucion}
  \textbf{Solución:} Multiplicamos por el recíproco de $\dfrac{2}{3}$, que es $\dfrac{3}{2}$:
  $$\frac{5}{6} \div \frac{2}{3} = \frac{5}{6} \cdot \frac{3}{2} = \frac{15}{12} = \frac{5}{4} = 1\frac{1}{4}.$$
\end{cuadrosolucion}

\textbf{Ejemplo R3.} Calcula: $a) \; \dfrac{3}{4} \cdot 2\frac{1}{2} \qquad b) \; 6 \div \dfrac{2}{5}$

\begin{cuadrosolucion}
  \textbf{Solución:}
  \begin{itemize}
    \item $a) \; 2\frac{1}{2} = \dfrac{5}{2}$, entonces $\dfrac{3}{4} \cdot \dfrac{5}{2} = \dfrac{15}{8} = 1\frac{7}{8}$.
    \item $b) \; 6 \div \dfrac{2}{5} = 6 \cdot \dfrac{5}{2} = \dfrac{6}{1} \cdot \dfrac{5}{2} = \dfrac{30}{2} = 15$.
  \end{itemize}
\end{cuadrosolucion}

\textbf{Ejemplo R4.} Calcula: $\dfrac{2}{3} \div \dfrac{5}{6} \cdot \dfrac{3}{4}$

\begin{cuadrosolucion}
  \textbf{Solución:} Resolvemos en orden de izquierda a derecha:
  $$\frac{2}{3} \div \frac{5}{6} \cdot \frac{3}{4} = \frac{2}{3} \cdot \frac{6}{5} \cdot \frac{3}{4} = \frac{2 \cdot 6 \cdot 3}{3 \cdot 5 \cdot 4} = \frac{36}{60} = \frac{3}{5}.$$
\end{cuadrosolucion}

%\newpage %ajuste pag

\subsection{Ejercicios propuestos}

\begin{enumerate}[leftmargin=*, resume]
  \item \facil{} Calcula: $a) \; \dfrac{2}{5} \cdot \dfrac{3}{7} \qquad b) \; \dfrac{1}{4} \cdot \dfrac{8}{9} \qquad c) \; \dfrac{5}{6} \cdot \dfrac{3}{10}$

  \item \facil{} Calcula: $a) \; \dfrac{3}{4} \div \dfrac{2}{5} \qquad b) \; \dfrac{7}{8} \div \dfrac{7}{8} \qquad c) \; \dfrac{9}{10} \div \dfrac{3}{5}$

  \item \facil{} Calcula: $a) \; 5 \cdot \dfrac{2}{3} \qquad b) \; \dfrac{3}{4} \cdot 8 \qquad c) \; \dfrac{1}{2} \div 4$

  \item \medio{} Calcula: $a) \; 2\frac{1}{3} \cdot 1\frac{1}{2} \qquad b) \; 3\frac{3}{4} \div \dfrac{5}{2}$

  \item \medio{} Calcula: $a) \; \dfrac{3}{5} \cdot \dfrac{10}{9} \cdot \dfrac{3}{4} \qquad b) \; \dfrac{2}{3} \div \dfrac{4}{9} \div \dfrac{3}{2}$

  \item \medio{} Un terreno de $\dfrac{3}{4}$ de hectárea se divide en 6 parcelas iguales. ¿Qué fracción de hectárea mide cada parcela?

  \item \dificil{} Calcula: $\left(\dfrac{2}{3} \cdot \dfrac{9}{8}\right) \div \left(\dfrac{5}{6} \cdot \dfrac{3}{10}\right)$
\end{enumerate}

%\newpage

% ============================================================
\section{Operaciones Combinadas}
% ============================================================

\subsection{Jerarquía de las operaciones}

\begin{cuadroteoria}
  \textbf{Orden de resolución (jerarquía):}
  \begin{enumerate}[leftmargin=*]
    \item \textbf{Paréntesis} (y otros signos de agrupación: corchetes, llaves), resolviendo de adentro hacia afuera.
    \item \textbf{Potencias} y raíces.
    \item \textbf{Multiplicaciones} y divisiones, de izquierda a derecha.
    \item \textbf{Sumas} y restas, de izquierda a derecha.
  \end{enumerate}
\end{cuadroteoria}

\newpage %ajuste para pág

\begin{cuadropasos}
  \textbf{Estrategia para operaciones combinadas con fracciones:}
  \begin{enumerate}[leftmargin=*]
    \item Resuelve primero lo que está dentro de los signos de agrupación.
    \item Aplica potencias si las hay.
    \item Realiza multiplicaciones y divisiones.
    \item Suma y resta los resultados.
    \item Simplifica la fracción final.
  \end{enumerate}
\end{cuadropasos}

\subsection{Ejercicios resueltos}

\textbf{Ejemplo R1.} Calcula sin signos de agrupación: $\dfrac{1}{2} + \dfrac{2}{3} \cdot \dfrac{3}{4}$

\begin{cuadrosolucion}
  \textbf{Solución:} Primero la multiplicación:
  $$\frac{2}{3} \cdot \frac{3}{4} = \frac{6}{12} = \frac{1}{2}.$$
  Luego la suma:
  $$\frac{1}{2} + \frac{1}{2} = \frac{2}{2} = 1.$$
\end{cuadrosolucion}

\textbf{Ejemplo R2.} Calcula sin signos de agrupación: $\dfrac{3}{4} - \dfrac{5}{6} \div \dfrac{2}{3}$

\begin{cuadrosolucion}
  \textbf{Solución:} Primero la división:
  $$\frac{5}{6} \div \frac{2}{3} = \frac{5}{6} \cdot \frac{3}{2} = \frac{15}{12} = \frac{5}{4}.$$
  Luego la resta:
  $$\frac{3}{4} - \frac{5}{4} = \frac{3 - 5}{4} = \frac{-2}{4} = -\frac{1}{2}.$$
\end{cuadrosolucion}

\textbf{Ejemplo R3.} Calcula con signos de agrupación: $\left(\dfrac{1}{2} + \dfrac{1}{3}\right) \cdot \dfrac{6}{5}$

\begin{cuadrosolucion}
  \textbf{Solución:} Primero el paréntesis:
  $$\frac{1}{2} + \frac{1}{3} = \frac{3 + 2}{6} = \frac{5}{6}.$$
  Luego multiplicamos:
  $$\frac{5}{6} \cdot \frac{6}{5} = \frac{30}{30} = 1.$$
\end{cuadrosolucion}

\textbf{Ejemplo R4.} Calcula con signos de agrupación: $\dfrac{\frac{3}{4} + \frac{1}{2}}{\frac{5}{6} - \frac{1}{3}}$

\begin{cuadrosolucion}
  \textbf{Solución:} Resolvemos numerador y denominador por separado:
  \begin{itemize}
    \item Numerador: $\dfrac{3}{4} + \dfrac{1}{2} = \dfrac{3}{4} + \dfrac{2}{4} = \dfrac{5}{4}$.
    \item Denominador: $\dfrac{5}{6} - \dfrac{1}{3} = \dfrac{5}{6} - \dfrac{2}{6} = \dfrac{3}{6} = \dfrac{1}{2}$.
  \end{itemize}
  Luego:
  $$\frac{\frac{5}{4}}{\frac{1}{2}} = \frac{5}{4} \div \frac{1}{2} = \frac{5}{4} \cdot 2 = \frac{10}{4} = \frac{5}{2}.$$
\end{cuadrosolucion}

\textbf{Ejemplo R5.} Calcula con signos de agrupación: $\left[\left(\dfrac{1}{3} + \dfrac{1}{6}\right) \cdot 2 - \dfrac{1}{4}\right] \div \dfrac{5}{8}$

\begin{cuadrosolucion}
  \textbf{Solución:}
  \begin{enumerate}
    \item Paréntesis interior: $\dfrac{1}{3} + \dfrac{1}{6} = \dfrac{2}{6} + \dfrac{1}{6} = \dfrac{3}{6} = \dfrac{1}{2}$.
    \item Multiplicación: $\dfrac{1}{2} \cdot 2 = 1$.
    \item Resta: $1 - \dfrac{1}{4} = \dfrac{3}{4}$.
    \item División final: $\dfrac{3}{4} \div \dfrac{5}{8} = \dfrac{3}{4} \cdot \dfrac{8}{5} = \dfrac{24}{20} = \dfrac{6}{5}$.
  \end{enumerate}
\end{cuadrosolucion}

\subsection{Ejercicios propuestos}

\begin{enumerate}[leftmargin=*, resume]
  \item \facil{} Calcula sin signos de agrupación: $a) \; \dfrac{1}{2} \cdot \dfrac{4}{5} + \dfrac{1}{5} \qquad b) \; \dfrac{3}{4} - \dfrac{1}{2} \cdot \dfrac{1}{2}$

  \item \facil{} Calcula sin signos de agrupación: $a) \; \dfrac{5}{6} \div \dfrac{1}{3} - \dfrac{1}{2} \qquad b) \; \dfrac{2}{3} + \dfrac{3}{4} \div \dfrac{9}{8}$

  \item \medio{} Calcula: $a) \; \left(\dfrac{1}{4} + \dfrac{3}{4}\right) \cdot \dfrac{2}{5} \qquad b) \; \left(\dfrac{2}{3} - \dfrac{1}{6}\right) \div \dfrac{1}{2}$

  \item \medio{} Calcula: $a) \; \left(2\frac{1}{2} - \dfrac{3}{4}\right) \cdot \dfrac{4}{7} \qquad b) \; \dfrac{3}{5} \cdot \left(\dfrac{5}{6} + \dfrac{1}{3}\right)$

  \item \medio{} Calcula: $\dfrac{\frac{2}{3} + \frac{1}{6}}{\frac{5}{6} - \frac{1}{2}}$

  \item \medio{} Calcula: $\left[\left(\dfrac{1}{2} + \dfrac{1}{4}\right) \cdot \dfrac{4}{3} + \dfrac{1}{6}\right] \div \dfrac{5}{12}$

  \item \dificil{} Calcula: $\dfrac{\frac{3}{4} \cdot \frac{8}{9} - \frac{1}{3}}{\frac{5}{6} \div \frac{5}{12} + \frac{1}{2}}$

  \item \dificil{} Plantea y resuelve: tengo $\dfrac{3}{4}$ de litro de jugo. Bebo $\dfrac{1}{3}$ del total y luego regalo la mitad de lo que queda. ¿Cuánto jugo me queda al final?
\end{enumerate}

%\newpage %ajuste pag

% ============================================================
\section{Potenciación en $\mathbb{Q}$}
% ============================================================

\subsection{Definición y propiedades}

\begin{cuadroteoria}
  \textbf{Potencia:} $\left(\dfrac{a}{b}\right)^n = \underbrace{\dfrac{a}{b} \cdot \dfrac{a}{b} \cdots \dfrac{a}{b}}_{n \text{ veces}} = \dfrac{a^n}{b^n}$ \quad ($b \neq 0$).

  \textbf{Signo del resultado:}
  \begin{itemize}[leftmargin=*]
    \item Exponente \emph{par}: el resultado es positivo. Ejemplo: $\left(-\dfrac{2}{3}\right)^2 = \dfrac{4}{9}$.
    \item Exponente \emph{impar}: se conserva el signo de la base. Ejemplo: $\left(-\dfrac{2}{3}\right)^3 = -\dfrac{8}{27}$.
  \end{itemize}

  \textbf{Propiedades de las potencias} (con $a, b \neq 0$):
  \begin{center}
    \begin{tabular}{ll}
      \toprule
      \textbf{Propiedad}       & \textbf{Fórmula}                                 \\
      \midrule
      Producto de potencias    & $a^m \cdot a^n = a^{m+n}$                        \\
      Cociente de potencias    & $\dfrac{a^m}{a^n} = a^{m-n}$                     \\
      Potencia de una potencia & $(a^m)^n = a^{m \cdot n}$                        \\
      Potencia de un producto  & $(ab)^n = a^n \cdot b^n$                         \\
      Potencia de un cociente  & $\left(\dfrac{a}{b}\right)^n = \dfrac{a^n}{b^n}$ \\
      Exponente cero           & $a^0 = 1$                                        \\
      Exponente negativo       & $a^{-n} = \dfrac{1}{a^n}$                        \\
      \bottomrule
    \end{tabular}
  \end{center}
\end{cuadroteoria}

\begin{cuadroadvertencia}
  \textbf{Cuidado con los errores clásicos:}
  \begin{itemize}[leftmargin=*]
    \item $\left(\dfrac{a}{b}\right)^{-n} = \left(\dfrac{b}{a}\right)^{n}$: un exponente negativo \emph{invierte} la fracción (además de potenciar).
    \item $\dfrac{a^m}{a^n} \neq a^{m \cdot n}$: los exponentes se \emph{restan}.
    \item La potencia se aplica a toda la base: $\left(-\dfrac{2}{3}\right)^2 = \dfrac{4}{9}$, pero $-\dfrac{2^2}{3} = -\dfrac{4}{3}$ (el signo no está dentro de la base).
  \end{itemize}
\end{cuadroadvertencia}

\subsection{Ejercicios resueltos}

\textbf{Ejemplo R1.} Calcula: $a) \; \left(\dfrac{2}{3}\right)^2 \qquad b) \; \left(-\dfrac{3}{4}\right)^3 \qquad c) \; \left(-\dfrac{1}{2}\right)^4$

\begin{cuadrosolucion}
  \textbf{Solución:}
  \begin{itemize}
    \item $a) \; \left(\dfrac{2}{3}\right)^2 = \dfrac{2^2}{3^2} = \dfrac{4}{9}$.
    \item $b) \; \left(-\dfrac{3}{4}\right)^3 = \dfrac{(-3)^3}{4^3} = \dfrac{-27}{64} = -\dfrac{27}{64}$ (exponente impar, se conserva el signo).
    \item $c) \; \left(-\dfrac{1}{2}\right)^4 = \dfrac{(-1)^4}{2^4} = \dfrac{1}{16}$ (exponente par, resultado positivo).
  \end{itemize}
\end{cuadrosolucion}

\textbf{Ejemplo R2.} Calcula: $a) \; \left(\dfrac{2}{5}\right)^{-2} \qquad b) \; \left(-\dfrac{3}{2}\right)^{-3}$

\begin{cuadrosolucion}
  \textbf{Solución:} Un exponente negativo invierte la base y luego se potencia:
  \begin{itemize}
    \item $a) \; \left(\dfrac{2}{5}\right)^{-2} = \left(\dfrac{5}{2}\right)^2 = \dfrac{25}{4}$.
    \item $b) \; \left(-\dfrac{3}{2}\right)^{-3} = \left(-\dfrac{2}{3}\right)^3 = \dfrac{(-2)^3}{3^3} = -\dfrac{8}{27}$.
  \end{itemize}
\end{cuadrosolucion}

\textbf{Ejemplo R3.} Aplica las propiedades y simplifica: $\left(\dfrac{3}{5}\right)^2 \cdot \left(\dfrac{3}{5}\right)^3$

\begin{cuadrosolucion}
  \textbf{Solución:} Producto de potencias de igual base: se suman los exponentes.
  $$\left(\frac{3}{5}\right)^2 \cdot \left(\frac{3}{5}\right)^3 = \left(\frac{3}{5}\right)^{2+3} = \left(\frac{3}{5}\right)^5 = \frac{243}{3125}.$$
\end{cuadrosolucion}

\newpage %ajuste pag

\textbf{Ejemplo R4.} Simplifica: $\left(\dfrac{81a^3b^2}{27a^5b^6}\right)^{-2}$

\begin{cuadrosolucion}
  \textbf{Solución:}
  \begin{enumerate}
    \item Simplificamos la fracción interior: $\dfrac{81}{27} = 3$, y aplicamos cociente de potencias a cada variable:
          $$\frac{81a^3b^2}{27a^5b^6} = 3 \cdot a^{3-5} \cdot b^{2-6} = 3a^{-2}b^{-4}.$$
    \item Aplicamos el exponente $-2$ (potencia de un producto y de una potencia):
          $$\left(3a^{-2}b^{-4}\right)^{-2} = 3^{-2} \cdot a^{(-2)(-2)} \cdot b^{(-4)(-2)} = 3^{-2}a^{4}b^{8}.$$
    \item Expresamos con exponentes positivos: $3^{-2} = \dfrac{1}{3^2} = \dfrac{1}{9}$, entonces:
          $$\left(\frac{81a^3b^2}{27a^5b^6}\right)^{-2} = \frac{a^4b^8}{9}.$$
  \end{enumerate}
\end{cuadrosolucion}

\textbf{Ejemplo R5.} Simplifica: $\left(\dfrac{2x^2y}{3z}\right)^3 \cdot \left(\dfrac{9z^2}{4x^4y^2}\right)$

\begin{cuadrosolucion}
  \textbf{Solución:}
  \begin{enumerate}
    \item Potencia del primer paréntesis: $\left(\dfrac{2x^2y}{3z}\right)^3 = \dfrac{8x^6y^3}{27z^3}$.
    \item Multiplicamos: $\dfrac{8x^6y^3}{27z^3} \cdot \dfrac{9z^2}{4x^4y^2} = \dfrac{8 \cdot 9}{27 \cdot 4} \cdot x^{6-4} \cdot y^{3-2} \cdot z^{2-3}$.
    \item Simplificamos los números: $\dfrac{72}{108} = \dfrac{2}{3}$. Entonces:
          $$\left(\frac{2x^2y}{3z}\right)^3 \cdot \left(\frac{9z^2}{4x^4y^2}\right) = \frac{2}{3}x^2y\, z^{-1} = \frac{2x^2y}{3z}.$$
  \end{enumerate}
\end{cuadrosolucion}

%\newpage

\subsection{Ejercicios propuestos}

\begin{enumerate}[leftmargin=*, resume]
  \item \facil{} Calcula: $a) \; \left(\dfrac{3}{4}\right)^2 \qquad b) \; \left(\dfrac{2}{5}\right)^3 \qquad c) \; \left(-\dfrac{1}{3}\right)^2$

  \item \facil{} Calcula: $a) \; \left(-\dfrac{2}{3}\right)^3 \qquad b) \; \left(\dfrac{4}{5}\right)^0 \qquad c) \; \left(-\dfrac{1}{2}\right)^5$

  \item \facil{} Calcula: $a) \; \left(\dfrac{3}{5}\right)^{-1} \qquad b) \; \left(\dfrac{2}{3}\right)^{-2} \qquad c) \; \left(-\dfrac{5}{4}\right)^{-2}$

  \item \medio{} Aplica propiedades: $a) \; \left(\dfrac{2}{7}\right)^3 \cdot \left(\dfrac{2}{7}\right)^4 \qquad b) \; \dfrac{\left(\frac{3}{5}\right)^7}{\left(\frac{3}{5}\right)^4}$

  \item \medio{} Simplifica: $a) \; \left(\dfrac{4x^3}{y^2}\right)^2 \qquad b) \; \left(\dfrac{2a^2b}{3c}\right)^3$

  \item \medio{} Simplifica: $\left(\dfrac{16x^4y^3}{8x^6y^7}\right)^{-1}$

  \item \medio{} Simplifica: $\left(\dfrac{27a^5b^2}{9a^2b^6}\right)^{-2}$

  \item \dificil{} Simplifica: $\left(\dfrac{2x^2y^3}{z^2}\right)^{-3} \cdot \left(\dfrac{x^4z^6}{8y^9}\right)$

  \item \dificil{} Simplifica: $\left(\dfrac{a^{-2}b^3}{a^4b^{-1}}\right)^{-2} \cdot \left(\dfrac{a^3b^{-2}}{b^5}\right)$

  \item \dificil{} El volumen de un cubo es $\left(\dfrac{2}{3}\right)^3\ \text{m}^3$. ¿Cuál es la medida de su arista?
\end{enumerate}

%\newpage

% ============================================================
\section{Ecuaciones Lineales con Fracciones}
% ============================================================

\subsection{¿Cómo se resuelven?}

\begin{cuadroteoria}
  \textbf{Ecuación lineal con fracciones:} es una ecuación de la forma $ax + b = 0$ (con $a \neq 0$) donde aparecen fracciones. Se resuelve con las mismas propiedades de la igualdad, pero conviene \emph{eliminar los denominadores} multiplicando toda la ecuación por el mcm de los denominadores.
\end{cuadroteoria}

\begin{cuadropasos}
  \textbf{Método de resolución:}
  \begin{enumerate}[leftmargin=*]
    \item \textbf{Eliminar denominadores:} multiplica \emph{toda} la ecuación por el mcm de los denominadores.
    \item \textbf{Distribuir:} aplica la propiedad distributiva en los paréntesis que queden.
    \item \textbf{Reducir:} agrupa los términos con $x$ en un lado y los números en el otro.
    \item \textbf{Despejar:} divide por el coeficiente de $x$.
    \item \textbf{Comprobar:} sustituye la solución en la ecuación original.
  \end{enumerate}
\end{cuadropasos}

\newpage %ajuste pag

\begin{cuadroadvertencia}
  \textbf{Errores clásicos:}
  \begin{itemize}[leftmargin=*]
    \item Multiplicar solo un término por el mcm en lugar de \emph{toda} la ecuación.
    \item Olvidar que el mcm debe multiplicar también los términos enteros.
    \item No simplificar la fracción final de la solución.
    \item Saltarse la comprobación: siempre verifica tu respuesta.
  \end{itemize}
\end{cuadroadvertencia}

\subsection{Ejercicios resueltos}

\textbf{Ejemplo R1.} Resuelve $\dfrac{x}{3} + \dfrac{x}{2} = 5$

\begin{cuadrosolucion}
  \textbf{Solución:}
  \begin{enumerate}
    \item mcm$(3, 2) = 6$. Multiplicamos toda la ecuación por 6:
          $$\cancel{6} \cdot \frac{x}{\cancel{3}} + \cancel{6} \cdot \frac{x}{\cancel{2}} = 6 \cdot 5 \quad \Rightarrow \quad 2x + 3x = 30.$$
    \item Reducimos: $5x = 30$.
    \item Despejamos: $x = \dfrac{30}{5} = 6$.
    \item Comprobación: $\dfrac{6}{3} + \dfrac{6}{2} = 2 + 3 = 5$ $\checkmark$
  \end{enumerate}
\end{cuadrosolucion}

\textbf{Ejemplo R2.} Resuelve $\dfrac{2x - 1}{3} = \dfrac{x + 2}{2}$

\begin{cuadrosolucion}
  \textbf{Solución:}
  \begin{enumerate}
    \item mcm$(3, 2) = 6$:
          $$\cancel{6} \cdot \frac{2x - 1}{\cancel{3}} = \cancel{6} \cdot \frac{x + 2}{\cancel{2}} \quad \Rightarrow \quad 2(2x - 1) = 3(x + 2).$$
    \item Distribuimos: $4x - 2 = 3x + 6$.
    \item $4x - 3x = 6 + 2 \Rightarrow x = 8$.
    \item Comprobación: $\dfrac{16 - 1}{3} = 5$ y $\dfrac{8 + 2}{2} = 5$ $\checkmark$
  \end{enumerate}
\end{cuadrosolucion}

\newpage

\textbf{Ejemplo R3.} Resuelve $\dfrac{3x}{4} - \dfrac{x}{2} = \dfrac{1}{6}$

\begin{cuadrosolucion}
  \textbf{Solución:}
  \begin{enumerate}
    \item mcm$(4, 2, 6) = 12$:
          $$12 \cdot \frac{3x}{4} - 12 \cdot \frac{x}{2} = 12 \cdot \frac{1}{6} \quad \Rightarrow \quad 9x - 6x = 2.$$
    \item Reducimos: $3x = 2 \Rightarrow x = \dfrac{2}{3}$.
    \item Comprobación: $\dfrac{3 \cdot \frac{2}{3}}{4} - \dfrac{\frac{2}{3}}{2} = \dfrac{2}{4} - \dfrac{2}{6} = \dfrac{1}{2} - \dfrac{1}{3} = \dfrac{1}{6}$ $\checkmark$
  \end{enumerate}
\end{cuadrosolucion}

\textbf{Ejemplo R4.} Resuelve $\dfrac{x + 3}{4} - \dfrac{x - 1}{6} = 2$

\begin{cuadrosolucion}
  \textbf{Solución:}
  \begin{enumerate}
    \item mcm$(4, 6) = 12$:
          $$\cancel{12} \cdot \frac{x + 3}{\cancel{4}} - \cancel{12} \cdot \frac{x - 1}{\cancel{6}} = 12 \cdot 2 \quad \Rightarrow \quad 3(x + 3) - 2(x - 1) = 24.$$
    \item Distribuimos con cuidado (el signo menos afecta a todo el segundo paréntesis):
          $$3x + 9 - 2x + 2 = 24.$$
    \item Reducimos: $x + 11 = 24 \Rightarrow x = 13$.
    \item Comprobación: $\dfrac{13 + 3}{4} - \dfrac{13 - 1}{6} = \dfrac{16}{4} - \dfrac{12}{6} = 4 - 2 = 2$ $\checkmark$
  \end{enumerate}
\end{cuadrosolucion}

\textbf{Ejemplo R5.} Problema: la mitad de un número más su tercera parte es igual a 25. ¿Cuál es el número?

\begin{cuadrosolucion}
  \textbf{Solución:} Sea $x$ el número buscado.
  $$\frac{x}{2} + \frac{x}{3} = 25.$$
  mcm$(2, 3) = 6$:
  $$3x + 2x = 150 \Rightarrow 5x = 150 \Rightarrow x = 30.$$
  Comprobación: $\dfrac{30}{2} + \dfrac{30}{3} = 15 + 10 = 25$ $\checkmark$ \quad El número es $30$.
\end{cuadrosolucion}

\newpage

\subsection{Ejercicios propuestos}

\begin{enumerate}[leftmargin=*, resume]
  \item \facil{} Resuelve: $\dfrac{x}{5} = 3$

  \item \facil{} Resuelve: $\dfrac{x}{4} + 2 = 5$

  \item \facil{} Resuelve: $\dfrac{x}{2} + \dfrac{x}{3} = 10$

  \item \facil{} Resuelve: $\dfrac{2x}{5} = \dfrac{6}{5}$

  \item \medio{} Resuelve: $\dfrac{x + 1}{3} = \dfrac{x - 2}{4}$

  \item \medio{} Resuelve: $\dfrac{3x}{4} - \dfrac{x}{8} = 5$

  \item \medio{} Resuelve: $\dfrac{2x - 3}{5} + \dfrac{x}{2} = 3$

  \item \medio{} Resuelve: $\dfrac{x}{3} + \dfrac{x + 1}{2} = \dfrac{7}{6}$

  \item \medio{} Resuelve: $\dfrac{2x + 1}{3} - \dfrac{x - 2}{4} = \dfrac{5x}{6}$

  \item \dificil{} Resuelve: $\dfrac{x - 1}{2} - \dfrac{x - 2}{3} = \dfrac{x - 3}{4}$

  \item \dificil{} La cuarta parte de un número aumentada en $\dfrac{1}{2}$ es igual a la sexta parte del número más 2. ¿Cuál es el número?

  \item \dificil{} En una clase, $\dfrac{2}{5}$ de los estudiantes practican fútbol, $\dfrac{1}{3}$ practican básquetbol y los 8 restantes practican otro deporte. ¿Cuántos estudiantes tiene la clase?
\end{enumerate}

\newpage

% ============================================================
\ifmwclaves
\section{Clave de Respuestas}
% ============================================================

\subsection*{Sección 1: Fracciones (Ejercicios 1--8)}

\begin{enumerate}[leftmargin=*, label=\textbf{\arabic*.}]
  \item a) propia \quad b) impropia \quad c) impropia ($=1$) \quad d) propia

  \item a) $\dfrac{11}{4}$ \quad b) $\dfrac{31}{6}$ \quad c) $\dfrac{26}{3}$

  \item a) $3\frac{2}{5}$ \quad b) $3\frac{1}{7}$ \quad c) $7\frac{3}{4}$

  \item a) $\dfrac{2}{3}$ \quad b) $\dfrac{2}{3}$ \quad c) $\dfrac{2}{3}$

  \item a) $6$ ($\frac{6}{15}$) \quad b) $28$ ($\frac{28}{32}$) \quad c) $24$ ($\frac{20}{24}$)

  \item Quedan $\dfrac{5}{8}$ de la pizza; es propia.

  \item $\dfrac{3}{4} \cdot 2 = \dfrac{6}{4} = \dfrac{3}{2} = 1\frac{1}{2}$ km.

  \item mcm$(4,6,12) = 12$: $\dfrac{9}{12}, \dfrac{10}{12}, \dfrac{7}{12}$ $\Rightarrow$ $\dfrac{7}{12} < \dfrac{3}{4} < \dfrac{5}{6}$.
\end{enumerate}

\newpage

\subsection*{Sección 2: Suma y Resta (Ejercicios 9--15)}

\begin{enumerate}[leftmargin=*, label=\textbf{\arabic*.}]
  \setcounter{enumi}{8}
  \item a) $\dfrac{1}{2}$ \quad b) $\dfrac{1}{2}$ \quad c) $1$

  \item a) $\dfrac{5}{6}$ \quad b) $\dfrac{3}{20}$ \quad c) $\dfrac{29}{35}$

  \item a) $\dfrac{11}{12}$ \quad b) $\dfrac{7}{24}$ \quad c) $\dfrac{11}{10} = 1\frac{1}{10}$

  \item a) $1\frac{1}{3} = \frac{4}{3}$; $\dfrac{4}{3} + \dfrac{5}{6} = \dfrac{8}{6} + \dfrac{5}{6} = \dfrac{13}{6} = 2\frac{1}{6}$ \quad b) $3\frac{1}{4} - 1\frac{1}{2} = \dfrac{13}{4} - \dfrac{3}{2} = \dfrac{13}{4} - \dfrac{6}{4} = \dfrac{7}{4} = 1\frac{3}{4}$

  \item a) $\dfrac{9}{12} + \dfrac{10}{12} - \dfrac{4}{12} = \dfrac{15}{12} = \dfrac{5}{4}$ \quad b) $\dfrac{7}{12} - \dfrac{3}{12} + \dfrac{8}{12} = \dfrac{12}{12} = 1$

  \item $\dfrac{2}{3} + \dfrac{1}{4} = \dfrac{8}{12} + \dfrac{3}{12} = \dfrac{11}{12}$ de taza.

  \item $1 - \left(\dfrac{1}{3} + \dfrac{1}{4}\right) = 1 - \dfrac{7}{12} = \dfrac{5}{12}$ del total.
\end{enumerate}

\newpage

\subsection*{Sección 3: Multiplicación y División (Ejercicios 16--22)}

\begin{enumerate}[leftmargin=*, label=\textbf{\arabic*.}]
  \setcounter{enumi}{15}
  \item a) $\dfrac{6}{35}$ \quad b) $\dfrac{8}{36} = \dfrac{2}{9}$ \quad c) $\dfrac{15}{60} = \dfrac{1}{4}$

  \item a) $\dfrac{3}{4} \cdot \dfrac{5}{2} = \dfrac{15}{8} = 1\frac{7}{8}$ \quad b) $1$ \quad c) $\dfrac{9}{10} \cdot \dfrac{5}{3} = \dfrac{45}{30} = \dfrac{3}{2}$

  \item a) $\dfrac{10}{3} = 3\frac{1}{3}$ \quad b) $6$ \quad c) $\dfrac{1}{8}$

  \item a) $\dfrac{7}{3} \cdot \dfrac{3}{2} = \dfrac{21}{6} = \dfrac{7}{2} = 3\frac{1}{2}$ \quad b) $\dfrac{15}{4} \cdot \dfrac{2}{5} = \dfrac{30}{20} = \dfrac{3}{2}$

  \item a) $\dfrac{3 \cdot 10 \cdot 3}{5 \cdot 9 \cdot 4} = \dfrac{90}{180} = \dfrac{1}{2}$ \quad b) $\dfrac{2}{3} \cdot \dfrac{9}{4} \cdot \dfrac{2}{3} = \dfrac{36}{36} = 1$

  \item $\dfrac{3}{4} \div 6 = \dfrac{3}{4} \cdot \dfrac{1}{6} = \dfrac{3}{24} = \dfrac{1}{8}$ de hectárea.

  \item $\left(\dfrac{2}{3} \cdot \dfrac{9}{8}\right) \div \left(\dfrac{5}{6} \cdot \dfrac{3}{10}\right) = \dfrac{18}{24} \div \dfrac{15}{60} = \dfrac{3}{4} \div \dfrac{1}{4} = 3$
\end{enumerate}

\newpage

\subsection*{Sección 4: Operaciones Combinadas (Ejercicios 23--30)}

\begin{enumerate}[leftmargin=*, label=\textbf{\arabic*.}]
  \setcounter{enumi}{22}
  \item a) $\dfrac{1}{2} \cdot \dfrac{4}{5} + \dfrac{1}{5} = \dfrac{2}{5} + \dfrac{1}{5} = \dfrac{3}{5}$ \quad b) $\dfrac{3}{4} - \dfrac{1}{4} = \dfrac{1}{2}$

  \item a) $\dfrac{5}{6} \cdot 3 - \dfrac{1}{2} = \dfrac{15}{6} - \dfrac{1}{2} = \dfrac{5}{2} - \dfrac{1}{2} = 2$ \quad b) $\dfrac{2}{3} + \dfrac{3}{4} \cdot \dfrac{8}{9} = \dfrac{2}{3} + \dfrac{2}{3} = \dfrac{4}{3}$

  \item a) $1 \cdot \dfrac{2}{5} = \dfrac{2}{5}$ \quad b) $\left(\dfrac{4}{6} - \dfrac{1}{6}\right) \div \dfrac{1}{2} = \dfrac{3}{6} \cdot 2 = 1$

  \item a) $\left(\dfrac{5}{2} - \dfrac{3}{4}\right) \cdot \dfrac{4}{7} = \dfrac{7}{4} \cdot \dfrac{4}{7} = 1$ \quad b) $\dfrac{3}{5} \cdot \dfrac{7}{6} = \dfrac{21}{30} = \dfrac{7}{10}$

  \item $\dfrac{\frac{5}{6}}{\frac{1}{3}} = \dfrac{5}{6} \cdot 3 = \dfrac{15}{6} = \dfrac{5}{2}$

  \item $\left[\left(\dfrac{3}{4}\right) \cdot \dfrac{4}{3} + \dfrac{1}{6}\right] \div \dfrac{5}{12} = \left(1 + \dfrac{1}{6}\right) \div \dfrac{5}{12} = \dfrac{7}{6} \cdot \dfrac{12}{5} = \dfrac{14}{5}$

  \item Numerador: $\dfrac{3}{4} \cdot \dfrac{8}{9} - \dfrac{1}{3} = \dfrac{2}{3} - \dfrac{1}{3} = \dfrac{1}{3}$. Denominador: $\dfrac{5}{6} \cdot \dfrac{12}{5} + \dfrac{1}{2} = 2 + \dfrac{1}{2} = \dfrac{5}{2}$. Resultado: $\dfrac{1}{3} \div \dfrac{5}{2} = \dfrac{2}{15}$.

  \item Bebo $\dfrac{1}{3} \cdot \dfrac{3}{4} = \dfrac{1}{4}$ L; quedan $\dfrac{1}{2}$ L; regalo $\dfrac{1}{4}$ L; me quedan $\dfrac{1}{4}$ L.
\end{enumerate}

\newpage

\subsection*{Sección 5: Potenciación (Ejercicios 31--40)}

\begin{enumerate}[leftmargin=*, label=\textbf{\arabic*.}]
  \setcounter{enumi}{30}
  \item a) $\dfrac{9}{16}$ \quad b) $\dfrac{8}{125}$ \quad c) $\dfrac{1}{9}$

  \item a) $-\dfrac{8}{27}$ \quad b) $1$ \quad c) $-\dfrac{1}{32}$

  \item a) $\dfrac{5}{3}$ \quad b) $\dfrac{9}{4}$ \quad c) $\dfrac{16}{25}$

  \item a) $\left(\dfrac{2}{7}\right)^7 = \dfrac{128}{823543}$ \quad b) $\left(\dfrac{3}{5}\right)^3 = \dfrac{27}{125}$

  \item a) $\dfrac{16x^6}{y^4}$ \quad b) $\dfrac{8a^6b^3}{27c^3}$

  \item $\left(\dfrac{16x^4y^3}{8x^6y^7}\right)^{-1} = \left(2x^{-2}y^{-4}\right)^{-1} = 2^{-1}x^2y^4 = \dfrac{x^2y^4}{2}$

  \item $\left(\dfrac{27a^5b^2}{9a^2b^6}\right)^{-2} = \left(3a^3b^{-4}\right)^{-2} = 3^{-2}a^{-6}b^8 = \dfrac{b^8}{9a^6}$

  \item $\left(\dfrac{2x^2y^3}{z^2}\right)^{-3} = \dfrac{z^6}{8x^6y^9}$; multiplicando: $\dfrac{z^6}{8x^6y^9} \cdot \dfrac{x^4z^6}{8y^9} = \dfrac{z^{12}}{64x^2y^{18}}$

  \item $\dfrac{a^{-2}b^3}{a^4b^{-1}} = a^{-6}b^4$; $\left(a^{-6}b^4\right)^{-2} = a^{12}b^{-8}$. Y $\dfrac{a^3b^{-2}}{b^5} = a^3b^{-7}$. Producto: $a^{15}b^{-15} = \dfrac{a^{15}}{b^{15}} = \left(\dfrac{a}{b}\right)^{15}$

  \item $V = \dfrac{8}{27}\ \text{m}^3 \Rightarrow a = \sqrt[3]{\dfrac{8}{27}} = \dfrac{2}{3}$ m.
\end{enumerate}

\newpage

\subsection*{Sección 6: Ecuaciones con Fracciones (Ejercicios 41--52)}

\begin{enumerate}[leftmargin=*, label=\textbf{\arabic*.}]
  \setcounter{enumi}{40}
  \item $x = 15$

  \item $\dfrac{x}{4} = 3 \Rightarrow x = 12$

  \item $3x + 2x = 60 \Rightarrow x = 12$

  \item $x = 3$

  \item $4(x + 1) = 3(x - 2) \Rightarrow 4x + 4 = 3x - 6 \Rightarrow x = -10$

  \item $6x - x = 40 \Rightarrow 5x = 40 \Rightarrow x = 8$

  \item $\dfrac{4x - 6}{10} + \dfrac{5x}{10} = 3 \Rightarrow 9x - 6 = 30 \Rightarrow x = 4$

  \item mcm $= 6$: $2x + 3(x + 1) = 7 \Rightarrow 5x = 4 \Rightarrow x = \dfrac{4}{5}$

  \item mcm $= 12$: $4(2x + 1) - 3(x - 2) = 10x \Rightarrow 8x + 4 - 3x + 6 = 10x \Rightarrow 10 = 5x \Rightarrow x = 2$

  \item mcm $= 12$: $6(x - 1) - 4(x - 2) = 3(x - 3) \Rightarrow 6x - 6 - 4x + 8 = 3x - 9 \Rightarrow 2x + 2 = 3x - 9 \Rightarrow x = 11$

  \item $\dfrac{x}{4} + \dfrac{1}{2} = \dfrac{x}{6} + 2$. mcm $= 12$: $3x + 6 = 2x + 24 \Rightarrow x = 18$

  \item $\dfrac{2x}{5} + \dfrac{x}{3} + 8 = x$. mcm $= 15$: $6x + 5x + 120 = 15x \Rightarrow 120 = 4x \Rightarrow x = 30$ estudiantes.
\end{enumerate}

\fi
\end{document}
